\scanpage{102}

\mspage{VII - 7 -}

n\textsuperscript{o}339

3) Soient $F$ un corps, et $R_1,\ldots,R_n$ des sous-anneaux de $F$ qui soient des anneaux de Krull. Alors leur \uline{intersection} $S=\handwritten{\bigcap}_{j=1}^{n}R_j$ est un anneau de Krull.

En effet, pour $1\handwritten{\leq}j\handwritten{\leq}n$, soit $(v_{ji})_{i\handwritten{\in}I_j}$ une famille de valuations du corps des fractions de $R_j$ vérifiant (AK.1), (AK.2), (AK.3) (où $A$ est remplacé par $R_j$). Notons $w_{ji}$ la restriction de $v_{ji}$ au corps des fractions de $S$. Alors la famille $(v_{ji})_{1\handwritten{\leq}j\handwritten{\leq}n,\ i\handwritten{\in}I_j}$ vérifie évidemment (AK.2) (où $A$ est remplacé par $S$), et aussi (AK.3) puisque l'ensemble des indices $j$ est fini. Les valuations $w_{ji}$ sont, soit discrètes, soit impropres. En ne gardant que celles qui sont discrètes, on obtient évidemment une famille vérifiant (AK.1), (AK.2) et (AK.3) (où $A$ est remplacé par $S$). Donc $S$ est bien un anneau de Krull.

4) En particulier, si $A$ est un anneau de Krull et $K'$ un sous-corps du corps des fractions $K$ de $A$, $K'\handwritten{\cap}A$ est un anneau de Krull.

THEOREME 1.- \uline{Soit} $A$ \uline{un anneau \unsure{intégré|intégre|intègre}}. \uline{Pour que} $A$ \uline{soit un anneau de Krull il faut et il suffit que les deux conditions suivantes soient satisfaites} :

a) $A$ \uline{est complètement intégralement clos} ;

b) \uline{toute famille non vide d'idéaux entiers divisoriels de} $A$ \uline{admet un élément maximal}.

Soit $A$ un Anneau de Krull. Il est complètement intégralement clos (Chap.VI, \S4, n\textsuperscript{o}3, cor. de la prop.9). Soit $(v_i)_{i\handwritten{\in}I}$ une famille de valuations du corps des fractions $K$ de $A$ vérifiant (AK.1), (AK.2) et (AK.3). On peut supposer les $v_i$ normées. Pour tout $\underline a\handwritten{\in}I(A)$, nous poserons

\begingroup\csname tagsleft@true\endcsname
\[ v_i(\underline a)=\sup_{\underline a\handwritten{\subset}Ax}\,(v_i(x))\ ;\tag*{(2)} \]
\endgroup

on a $v_i(\underline a)\handwritten{\in}Z$, car, si $a$ est un élément non nul de $\underline a$, la relation $Ax\handwritten{\supset}Aa$ implique $v_i(x)\handwritten{\leq}v_i(a)$ (d'après (AK.2)), ce qui montre que la famille des $v_i(x)$ ($\underline a\handwritten{\subset}Ax$) est majorée. Etablissons les propriétés suivantes :

1) \uline{Soit} $\underline a$ \uline{un idéal divisoriel} ; \uline{pour que} $y\handwritten{\in}\underline a$, \uline{il faut et il suffit que}

% archive sheet number: 101
